% Complete base slice before imposing the mixed F equations.
\subsection{共点三直線における基底の非被約な形式切片}
\label{subsec:concurrent-base-v32}

$\mathscr Z=\{(c,D)\in S_3\oplus\mathfrak{sl}_3:D(c)=0\}$ と置く。
この基底スキーム自身に既に非被約構造が現れる。
以下の補題は混合項 $F$ の方程式をまだ課さない、基底だけの全形式的な記述である。

\begin{lemma}[共点三直線の基底切片]
\label{lem:concurrent-base-v32}
\[
 c_0=xy(x+y),\qquad D_0=(Ax+By)\partial_z
\]
とし、
\[
 U=B-2A,\qquad V=A-2B,\qquad W=A+B
\]
が全て非零とする。例えば $(A,B)=(2,5)$ はこの条件を満たす。
この点における $\mathscr Z$ の完備局所環は
\[
 \frac{\mathbb C[[u_1,\ldots,u_8,a,b]]}{(ab,b^2)}
\]
に同型である。八つの滑らかな変数のうち六つは $c_0$ の
$\mathrm{SL}_3$ 軌道方向、残る二つは $A,B$ の変動である。
被約化は滑らかな九次元であり、零でない平方零イデアル $(b)$ は
共点三直線の部分 $a=0$ に支持される。
\end{lemma}

\begin{proof}
$c_0$ の $\mathfrak{sl}_3$ による接軌道は六次元である。
その像は二変数三次式の四次元空間と
$z\partial_xc_0,z\partial_yc_0$ の二方向からなる。
従って、四次元の形式横断切片を
\[
 c=xy(x+y)+a xyz+b xz^2+\gamma yz^2+d z^3
\]
と取れる。六つの独立な軌道方向を形式的な群座標として選べば、
軌道方向とこの切片から $S_3$ への写像の接写像は同型であり、形式逆写像定理が適用できる。
$D$ も同じ $\mathrm{SL}_3$ 座標変換で移すので、$D(c)=0$ の問題は
この六次元の滑らかな因子と切片上の問題に分かれる。

切片上で $D_{zx}=A,D_{zy}=B$ と書く。
$D(c)=0$ の四つの二変数三次単項式と $x^2z,y^2z$ の係数から、
残る六つの係数は一意に
\[
\begin{aligned}
 D_{yx}=D_{xy}&=0,\\
 D_{xx}&=\frac{a(B-2A)}3,&
 D_{yy}&=\frac{a(A-2B)}3,\\
 D_{xz}&=-2B\gamma,&D_{yz}&=-2Ab,\\
 D_{zz}&=-D_{xx}-D_{yy}
\end{aligned}
\]
と決まる。代入後の残る四式は、零でない定数倍を除いて
\[
\begin{aligned}
 E_1&=Ub+V\gamma,\\
 E_2&=aUb+3Ad,\\
 E_3&=aV\gamma+3Bd,\\
 E_4&=W(ad-2b\gamma)
\end{aligned}
\]
である。

この段階で $U,V,W$ は形式局所環の単元である。まず
\[
 E_2+E_3-aE_1=3Wd
\]
から $d$ を消去でき、次に $E_2-3Ad=Uab$ から $ab=0$ を得る。
また $E_1=0$ により $\gamma=-Ub/V$ である。
最後に $E_4=0$ から $b\gamma=0$、従って $b^2=0$ が従う。
逆に $d=0,\gamma=-Ub/V,ab=b^2=0$ を満たせば四つの $E_i$ は全て零になる。
従ってこれは零点集合だけの比較でなく、局所化された関係イデアルの等式である。

残った自由変数は六つの軌道座標と $A,B,a,b$ であり、表示した完備局所環を得る。
$\mathbb C[[a,b]]/(ab,b^2)$ において $b\ne0$、$b^2=ab=0$ だから、
非被約性とその支持も明らかである。
なお、消去には $A$ や $B$ 自身の可逆性は必要なく、上記三つの単元条件で十分である。
\end{proof}

被約化の切片では $c=xy(x+y+az)$ となる。
従って $a\ne0$ は非共点の三直線、$a=0$ は共点三直線に対応する。
この補題だけから、混合項を含む全 Poisson スキームの既約成分性や
一般点の非被約厚みの次数を結論してはならない。
その判定には、この環上での混合写像をさらに解析する必要がある。
